\scanpage{88}

En \unsure{dernier|deuxième|troisième|anticipant} lieu, soit $f(\unsure{s|r|\rho})$ une fonction définie pour $s>0$, telle que $f(e^t)$ soit convexe, posons $f(s_1,\ldots,s_n)=\sum_1^n f(s_i)$, c'est une fonction symétrique des \deleted{\unsure{nombres}}\inserted{\unsure{argts|variables}} $s_i>0$, et telle que $f(e^{t_1},\ldots,e^{t_n})$ soit convexe. Par suite, \deleted{\unsure{$F$}} $f(s_1,\deleted{\unsure{+}}\ldots s_n)$ satisfait aux conditions du corollaire 1, sous l'hypothèse b. Si on suppose de plus $f(s)$ croissante, et définie et continue aussi pour $s\geq0$, alors $f(s_1,\ldots s_n)$ satisfait \unsure{même} aux conditions du \deleted{corollaire 1}\inserted{th. 1, 2°} sous l'hypothèse a.

D'où, \deleted{\unsure{avec la hypothèse}} \inserted{sous les conditions b,} \deleted{convenables sur $(\alpha_i),(\beta_i)$, \unsure{les inégalités}} l'inégalité

\[
\sum_1^n f(\rho_i)\leq\sum_1^n f(\sigma_i)\tag*{(13)}
\]

et sous les conditions de \unsure{b)|a)}, compte tenu d'une remarque précédente, la suite d'inégalités

\[
\sum_1^k f(\rho_i)\leq\sum_1^k f(\sigma_i)\qquad(k=1,\ldots n)\tag*{(14)}
\]

Prenons p.ex, \deleted{\unsure{$p$ positif et|pour}}\inserted{pour $p$ réel quelconque} $f(s)=s^p$, \inserted{$(s>0)$} on a alors $f(e^t)=e^{pt}$, qui est bien une fonction convexe. \deleted{\unsure{$f$ elle}}\inserted{$f(r)$} est définie et \deleted{\unsure{même}} continue pour $s\geq0$, et \deleted{\unsure{une}} fonction croissante de $s$, dès que $p\geq0$. \deleted{\unsure{On obtient donc}}

% archive sheet number: 89
